% !TEX root = main.tex
% APPENDIX D: the template's files. If you add or rename files, update the
% tree and the descriptions below (or delete this appendix from your thesis).

\chapter{Files in the template}
\label{app:filetree}

The template is a folder of \LaTeX\ source files. The tree below shows every file, grouped in the order \texttt{main.tex} reads them; all of them sit together in one folder except the images, which are in \texttt{Graphics/}. \cref{sec:filedescriptions} describes what each file contains. To build the thesis, compile \texttt{main.tex}; see \cref{sec:checklist} for what to check before you submit.

{\singlespacing\small
\begin{verbatim}
NJIT-Thesis-Template/
|-- readme.md                Overview and build instructions
|-- LICENSE                  License (LPPL) for the class file
|-- njit_thesis.cls          Class file: the NJIT formatting rules
|-- main.tex                 Master file; compile this one
|-- packages.tex             Packages and citation style (Appendix C)
|-- macros.tex               Your notation and theorem environments
|-- references.bib           Bibliography database
|
|-- abstract.tex             Front matter, in the order main.tex
|-- approval.tex               reads it; biography.tex also holds
|-- biography.tex              the title, name, degrees, and date
|-- dedication.tex
|-- acknowledgment.tex
|-- listofsymbols.tex        Two-column List of Symbols
|-- listofdefinitions.tex    Two-column Glossary of Terms
|
|-- chapter1.tex             Chapter 1: getting started
|-- chapter2.tex             Chapter 2: what the GSO requires
|-- table.tex                Sample table, \input by chapter2.tex
|-- chapter3.tex             Chapter 3: writing the text
|-- chapter4.tex             Chapter 4: figures and tables
|-- chapter5.tex             Chapter 5: references and submission
|-- appendixA.tex            Appendix A (several appendices)
|-- appendixB.tex            Appendix B: citation styles
|-- appendixC.tex            Appendix C: packages
|-- appendixD.tex            Appendix D: this list of files
|-- appendix.tex             Alternative: a single appendix
|-- appendix-rules.tex       Shared text, \input by appendix.tex
|-- appendix-demo.tex          and appendixA.tex
`-- Graphics/                Images for figures
    |-- FancyA.jpg
    |-- FancyB.png
    |-- EberhardtHall.jpg
    |-- drop_cap_1.png
    `-- drop_cap_2.png
\end{verbatim}
}

Compiling also creates \texttt{main.pdf} and a set of auxiliary files with the same base name (\texttt{main.aux}, \texttt{.bbl}, \texttt{.bcf}, \texttt{.blg}, \texttt{.fdb\_latexmk}, \texttt{.fls}, \texttt{.lof}, \texttt{.log}, \texttt{.lot}, \texttt{.run.xml}, \texttt{.synctex.gz}, and \texttt{.toc}). These are rebuilt on every run and can be deleted at any time.

\section{What Each File Contains}
\label{sec:filedescriptions}

Here is what each file contains and how to edit it.

\begin{description}
\emergencystretch=2.5em
  \item[\texttt{njit\_thesis.cls}] The \LaTeX\ class file. It encapsulates all official NJIT Graduate Studies Office formatting rules, including 1.5-inch left margins for binding, 1.0-inch top/bottom/right margins, double line spacing for body text, centered bold uppercase chapter titles, multi-page continuation headers for the Table of Contents, Lists of Figures and Tables, and correct appendix handling. You should not need to edit this class file directly.

  \item[\texttt{main.tex}] The master control document. It loads the class, includes packages and macros, sets the graphicspath, loads the front matter, and arranges the chapters and appendices via \verb+\input+ statements.

  \item[\texttt{packages.tex}] Loads optional \LaTeX\ packages (such as \texttt{tabularx}, \texttt{amsmath}, \texttt{amsthm}, the Times fonts \texttt{newtxtext} and \texttt{newtxmath}, \texttt{graphicx}, \texttt{subcaption}, \texttt{booktabs}, \texttt{xurl}, and \texttt{biblatex}). Add any packages your research needs there, before the \texttt{hyperref} line.

  \item[\texttt{macros.tex}] Defines custom mathematical shortcuts and theorem environments (e.g., Theorem, Lemma, Definition, Proposition, Corollary). The file contains some examples. Well-defined and consistently used macros can make writing and editing your thesis much easier.

  \item[\texttt{abstract.tex}] Enter the text of your abstract here. The class file automatically supplies the centered title and author header.

  \item[\texttt{approval.tex}] Defines the dissertation advisor and committee members using the \verb+\advisor+ and \verb+\member+ commands. If you have a co-advisor, add \verb+\coadvisor{Name}{Title, Institution}+ after \verb+\advisor+; both advisors are then labeled ``Co-Advisor'' and listed before the committee. Faculty titles must match the catalog.

  \item[\texttt{biography.tex}] Contains your dissertation title, author name, major, degree title, department, institution, graduation date (\verb+\graduationdate+), and degrees earned (entered with the \verb+\degreeearned+ command; the older name \verb+\degree+ still works unless a package such as \texttt{gensymb} defines it). Also includes the hook for \verb+\printmypublications+. Date and place of birth are added only in the copy submitted to the GSO.

  \item[\texttt{dedication.tex}] Contains your optional dedication. If you do not wish to include a dedication, simply comment out \verb+\dedication{% !TEX root = main.tex
% DEDICATION (optional; included via \dedication{% !TEX root = main.tex
% DEDICATION (optional; included via \dedication{% !TEX root = main.tex
% DEDICATION (optional; included via \dedication{\input{dedication}} in main.tex)
%
% A short dedication or quotation, set single-spaced in a centered block.
% Check the spelling of names and keep person and tense consistent (GSO).
% To omit the page, comment out the \dedication line in main.tex.
% Words left out of a quotation are marked with an ellipsis, typed \ldots;
% between two full sentences the first keeps its period (Chicago style).

\emph{Power of concentration is a very valuable quality, which few people acquire except through education. \ldots{} To be able to concentrate on the same matter for a considerable time is essential to difficult achievement, and even to the understanding of any complicated or abstruse subject.}

\begin{flushright}
Bertrand Russell, \emph{On Education} (1926)
\end{flushright}
} in main.tex)
%
% A short dedication or quotation, set single-spaced in a centered block.
% Check the spelling of names and keep person and tense consistent (GSO).
% To omit the page, comment out the \dedication line in main.tex.
% Words left out of a quotation are marked with an ellipsis, typed \ldots;
% between two full sentences the first keeps its period (Chicago style).

\emph{Power of concentration is a very valuable quality, which few people acquire except through education. \ldots{} To be able to concentrate on the same matter for a considerable time is essential to difficult achievement, and even to the understanding of any complicated or abstruse subject.}

\begin{flushright}
Bertrand Russell, \emph{On Education} (1926)
\end{flushright}
} in main.tex)
%
% A short dedication or quotation, set single-spaced in a centered block.
% Check the spelling of names and keep person and tense consistent (GSO).
% To omit the page, comment out the \dedication line in main.tex.
% Words left out of a quotation are marked with an ellipsis, typed \ldots;
% between two full sentences the first keeps its period (Chicago style).

\emph{Power of concentration is a very valuable quality, which few people acquire except through education. \ldots{} To be able to concentrate on the same matter for a considerable time is essential to difficult achievement, and even to the understanding of any complicated or abstruse subject.}

\begin{flushright}
Bertrand Russell, \emph{On Education} (1926)
\end{flushright}
}+ in \texttt{main.tex}.

  \item[\texttt{acknowledgment.tex}] Contains your acknowledgments. Thank your advisor first, then your committee; the GSO Front Matter guide then lists other contributors, family, funding sources, and technical support.

  \item[\texttt{listofsymbols.tex}] Contains a two-column table of symbols and their meanings, one \verb+\sym{symbol}{meaning}+ per row. Included via \verb+\symbols+ in \texttt{main.tex}.

  \item[\texttt{listofdefinitions.tex}] Contains a two-column table of terms and definitions in alphabetical order, one \verb+\term{term}{definition}+ per row. Included via \verb+\glossaryofterms+ in \texttt{main.tex}.


  \item[\texttt{chapter1.tex}, \texttt{chapter2.tex}, \ldots] The individual chapter files. Keeping chapters in separate files allows you to comment out finished chapters while drafting new ones to speed up compilation. The sample has five chapters: getting started, what the GSO requires, writing the text, figures and tables, and references and submission.

  \item[\texttt{table.tex}] The sample \cref{table:gso_rules}, kept in its own file and pulled into \texttt{chapter2.tex} with \verb+\input+, to show how a large or script-generated table can be kept out of the chapter text.

  \item[\texttt{appendixA.tex} to \texttt{appendixD.tex}] The sample's four appendices, each starting with \verb+\chapter+: how appendices work, guides and examples for the citation styles, the packages the template loads, and this list of files.

  \item[\texttt{appendix.tex}] The alternative for a document with a single appendix, which starts with \verb+\oneappendix+ instead; switch to it in \texttt{main.tex}. Sample text shared with \texttt{appendixA.tex} lives in \texttt{appendix-rules.tex} and \texttt{appendix-demo.tex}, pulled in with \verb+\input+; delete them when you write your own appendices.

  \item[\texttt{references.bib}] A standard \BibTeX\ database of citations. Any entry with the field \verb+keywords = {me}+ will also be included automatically in the Presentations and Publications section of your Biographical Sketch.

  \item[\texttt{Graphics/}] Directory containing all figures, plots, and diagrams. If an image that a figure uses is missing, such as \texttt{EberhardtHall.jpg} for \cref{fig:eberhardt}, the sample prints a framed placeholder instead.
\end{description}
